Condition on both pilot folds.  Independence of the evaluation observations from the opposite-fold pilots follows from the sampling assumption and the deterministic split.  Apply the exact Hoeffding decomposition to each diagonal-free quadratic correction.  After averaging folds, its first projection combines with the displayed one-step score to give \(\mathbb L_{J_y,J_z,n}\), its canonical second projection is \(\mathbb U_{J_y,J_z,n}\), and the omitted projection components are \(B^y_{J_y,n}+B^z_{J_z,n}\).  Thus the decomposition is an identity plus the pilot/operator remainder.

For the linear projection, the first variation in the cosine-second-variation lemma and uniform ellipticity show that the squared \(L^2(P_{0,n})\oplus L^2(P_{1,n})\) norm of the centered score is comparable to \(\rho_n^2\).  Independence of the two arms and the two equal folds therefore gives
\[
 \operatorname{Var}(\mathbb L_{J_y,J_z,n})\asymp\frac{\rho_n^2}{n}, \tag{1}
\]
with the score identically zero when \(\Delta_n=0\).

For the canonical quadratic projection, the two-sample feature contrast on one fold has covariance \(\Gamma/m\).  Whitening and diagonalizing \(K=\Gamma^{1/2}H\Gamma^{1/2}\), the exact variance formula for a diagonal-free order-two statistic gives
\[
 \operatorname{Var}(\mathbb U_{J_y,J_z,n})
 =\frac{4\sum_k\nu_{k,J_y,J_z,n}^2}{n^2}\{1+o(1)\}
 =\frac{A_{J_y,J_z,n}^2}{n^2}\{1+o(1)\}.                \tag{2}
\]
The \(o(1)\) is uniform because \(D=o(n^2)\), the densities are bounded, and compact wavelet support bounds the fourth-moment contraction by \(D/n^2=o(1)\).  The spectrum proposition supplies the two deterministic bias bounds.

The proposed pilot-control lemma bounds every remaining plug-in and ridge term by
\[
 O_P(\varepsilon_{\mathrm{pil},n}^3+
 \rho_n^{\,2/(2+3/s_y+d/s_z)}\varepsilon_{\mathrm{pil},n}^2)
 +o_P(1)\left\{\frac{\rho_n}{\sqrt n}+\frac A n\right\}. \tag{3}
\]
Under the proposed condition, (3) is little-oh of the sum of the four displayed scales.  Combining (1)--(3) with the exact Hoeffding identity proves the proposed uniform expansion.  At the rate-balancing resolutions, the spectrum proposition and the tuning identities give
\[
 \frac A n\asymp\frac{2^{dJ_z^*/2}}n
 \asymp2^{-2s_zJ_z^*}
 \asymp2^{-2(s_y+1)J_y^*}
 \asymp r_{n,d}(s_z).                                   \tag{4}
\]
Consequently \(\varepsilon_{\mathrm{pil},n}^3+\rho_n^{\,2/(2+3/s_y+d/s_z)}\varepsilon_{\mathrm{pil},n}^2=o\{\rho_n/\sqrt n+r_{n,d}(s_z)\}\) makes the observable remainder negligible at that scale.  Without this last condition the same calculation gives the honest root-mean-square upper scale \(\rho_n/\sqrt n+r_{n,d}(s_z)+\varepsilon_{\mathrm{pil},n}^3+\rho_n^{\,2/(2+3/s_y+d/s_z)}\varepsilon_{\mathrm{pil},n}^2\).